\PassOptionsToPackage{dvipsnames}{xcolor}
\documentclass[oneside,11pt]{book}
\usepackage{pagecolor}
\usepackage[utf8]{inputenc}
\usepackage[T1]{fontenc}
\usepackage{soul}
\usepackage[margin=1.5in]{geometry}
\usepackage{tcolorbox}
\tcbuselibrary{documentation, skins}
\usepackage{titlesec}

% \setsecnumdepth{subsection}
% \setcounter{tocdepth}{3}
\usepackage{fontawesome5}
\usepackage{censor}
\usepackage{enumitem}
\usepackage{epigraph}
\newcommand{\myepigraphsimpsons}[1]{
  \epigraph{\vspace{-0.2in} \emph{#1}}{\textsc{The Simpsons}}
}

\usepackage{cite}
\usepackage{caption}
\captionsetup{font=small}
\usepackage{graphicx}
\usepackage{pdfpages}
\usepackage{hyperref}
\usepackage{wrapfig}
\setlength\intextsep{0pt} % remove extra space above and below in-line float

\usepackage{booktabs}

\usepackage{tikz}
\usetikzlibrary{positioning,calc,arrows.meta,shapes.geometric}

%\usepackage[dvipsnames]{xcolor}
%\usepackage{anyfontsize}
%\usepackage{sectsty}
%\usepackage[acronym]{glossaries}
%\setacronymstyle{long-short-desc}
%\makeglossaries{}
%\usepackage[makeroom]{cancel}

% boxes
\newtcolorbox{marker}[1][]{enhanced,
  before skip balanced=2mm,after skip balanced=3mm,
  boxrule=0.4pt,left=5mm,right=2mm,top=1mm,bottom=1mm,
  colback=yellow!10,
  sharp corners,rounded corners=southeast,arc is angular,arc=3mm,
  underlay={%
    \path[fill=tcbcolback!80!black] ([yshift=3mm]interior.south east)--++(-0.4,-0.1)--++(0.1,-0.2);
    \path[draw=tcbcolframe,shorten <=-0.05mm,shorten >=-0.05mm] ([yshift=3mm]interior.south east)--++(-0.4,-0.1)--++(0.1,-0.2);
    \path[fill=tcbcolback!80!black,draw=none] (interior.south west) rectangle node{\large\bfseries {#1}} ([xshift=4mm]interior.north west);
    },
  drop fuzzy shadow}

\newenvironment{commentbox}[1][]{
  \small
  \begin{marker}[\faChessPawn{}]
    {\textbf{#1}}
  }{
  \end{marker}
}

\newenvironment{surprisebox}[1][]{
  \small
  \begin{marker}[\faChessKnight{}]
    {\textbf{#1}}
  }{
  \end{marker}
}

\newenvironment{examplebox}[1][]{
  \small
  \begin{marker}[\faChessBishop{}]
    {\textbf{#1}}
  }{
  \end{marker}
}

\newenvironment{questionbox}[1][]{
  \small
  \begin{marker}[\faChessRook{}]
    {\textbf{#1}}
  }{
  \end{marker}
}

\newenvironment{keybox}[1][]{
\small
  \begin{marker}[\faChessQueen{}]
    {\textbf{#1}}
  }{
  \end{marker}
}

\newenvironment{warningbox}[1][]{
  \small
  \begin{marker}[\faChessKing{}]
    {\textbf{#1}}
  }{
  \end{marker}
}

\renewcommand{\figurename}{Fig.}
\renewcommand{\tablename}{Tab.}
%\def\Section{\S}
\newcommand{\Section}{\textsection}

\renewcommand{\figureautorefname}{Fig.}
\renewcommand{\tableautorefname}{Tab.}
\makeatletter
\renewcommand{\chapterautorefname}{\textsection\@gobble}
\renewcommand{\sectionautorefname}{\textsection\@gobble}
\renewcommand{\subsectionautorefname}{\textsection\@gobble}
\renewcommand{\appendixautorefname}{\textsection\@gobble}

\newcommand*\l@chapterinfo{\@nodottedtocline{0}{0.0em}{1.5em}}
\newcommand*\l@sectioninfo{\@nodottedtocline{1}{1.5em}{2.3em}}
\newcommand*\l@subsectioninfo{\@nodottedtocline{2}{3.8em}{3.2em}}

% Copied from the book class macro `\@dottedtocline`. Removed the dots and page number
\def\@nodottedtocline#1#2#3#4#5{%
  \ifnum #1>\c@tocdepth \else
    \vskip \z@ \@plus.2\p@
    {\leftskip #2\relax \rightskip \@tocrmarg \parfillskip -\rightskip
     \parindent #2\relax\@afterindenttrue
     \interlinepenalty\@M
     \leavevmode
     \@tempdima #3\relax
     \advance\leftskip \@tempdima \null\nobreak\hskip -\leftskip
     {#4}\nobreak
     \leaders\hbox{$\m@th
        \mkern \@dotsep mu\hbox{\,}\mkern \@dotsep
        mu$}\hfill
     \nobreak
     \hb@xt@\@pnumwidth{\hfil\normalfont \normalcolor }%
     \par}%
  \fi}



\makeatother

\newcommand{\chapterinfo}[1]{%
  \phantomsection
  \addcontentsline{toc}{chapterinfo}{\textcolor{black}{\emph{#1}}}%
}
\newcommand{\sectioninfo}[1]{%
  \phantomsection
  \addcontentsline{toc}{sectioninfo}{\textcolor{black}{\emph{#1}}}%
}
\newcommand{\subsectioninfo}[1]{%
  \phantomsection
  \addcontentsline{toc}{subsectioninfo}{\textcolor{black}{\texttt #1}}%
}

\newcommand{\mycomment}[3][\color{blue}]{{#1{{#2}: {#3}}}}
\newcommand{\tvn}[1]{\mycomment{TVN}{#1}}{}
\newcommand{\didi}[1]{\mycomment{Didier}{#1}}{}
\newcommand{\tl}[1]{\mycomment{ThanhLe}{#1}}{}
\newcommand{\xz}[1]{\mycomment{Xiaokuan}{[#1]}}{}
\newcommand{\red}[1]{{\color{red}{#1}}}

\newcommand{\mytitle}{Roars Lab's Handbook}
%\newcommand{\mysubtitle}{A Handbook for Navigating CS PhD Admissions in the U.S.}
%\newcommand{\mybookaddress}{\href{https://roars.dev/phd-cs-us/demystify.pdf}{https://roars.dev/phd-cs-us/demystify.pdf}}
\newcommand{\mybookgithub}{https://code.roars.dev/phd-cs-us}
\newcommand{\me}{Roars' Members}
\newcommand{\institute}{George Mason University}

\usepackage{hyperref}
\hypersetup{
    hidelinks,
    linktoc=all,
  colorlinks,
  citecolor=black,
  filecolor=black,
  linkcolor=blue,
  urlcolor=blue,
  pdfpagemode=UseOutlines,
}





\begin{document}

\frontmatter
%\newpagecolor{SpringGreen}

\begin{titlepage}
    \begin{tikzpicture}[remember picture, overlay]
        \draw[OliveGreen, line width = 3pt] ($(current page.north west) + (1in,-1in)$) rectangle ($(current page.south east) + (-1in,1in)$);
      \end{tikzpicture}

    \centering
    {\Huge\bfseries\baselineskip1.3em\mytitle\par}\bigskip\bigskip
    \vfill
    \includegraphics[width=1\textwidth]{files/dino.png}
    \vfill
    {\Large\bfseries\me{}}\medskip\par
    {\large\institute{}}%\medskip\par    
    %{\small\today{}}\medskip\par
    %\mybookaddress{} (\href{\mybookgithub}{Github})
\end{titlepage}
%\restorepagecolor
\newpage
\pagestyle{empty}


% \begin{figure}[h]
%     \centering
%     \includegraphics[scale=0.1]{files/IMG_0785.jpeg}
%     \caption{Two of my angels. They are breeding and laying eggs.}\label{fig:angelfish}
%   \end{figure}
  


\renewcommand{\contentsname}{Contents and Summary}
\tableofcontents


\chapter{Preface}

\mainmatter
\part{Introduction}\label{part:intro}
\chapter{Meeting}

\begin{itemize}
    \item Meeting is mandatory for all lab members and is in-person.
    \item We meet every Thursday from 2pm to 3pm.  Typically, our meeting directly follows the Software Engineering seminar (Thurs from 1pm to 2pm).
    \item We meet in the main CS conference room (Room \#4201 Nguyen Engineering Building).
    \item If you need to attend remotely, you can join via Zoom but you need to inform me ahead of time.
\end{itemize}
\part{Writing}
\chapter{Papers}
\chapter{Rebuttal}

\begin{itemize}
\item Start with a \emph{thank you note} to the reviewers for their time and feedback. Reviewing is a thankless job and reviewers often put in lots of time and effort to provide feedback.  Even if the reviews are harsh and you disagree, do it anyway. It can make them more receptive to your points.
\end{itemize}
\chapter{Presentation}

\begin{itemize}
\item \textbf{Number your slides}.  This is useful when people want to ask questions (or give feedback) on specific slides (e.g., \emph{``on slide 10, you said ...''}).
\item \textbf{Every slide takes approximately 1 minutes}. Using this rule of thumb, you can estimate how many slides you need for your presentation.
\end{itemize}

\part{Classes}

\begin{itemize}
\item Independent Study
\begin{itemize}
    \item You need to be my student before you can register for independent study with me.  I.e., unless I officially advise you, you cannot register for independent study with me.
    \item You will typically get an A grade if you do the work satisfactorily. This is one of the reason why PhD students often get high GPAs, and also the reason why no one cares about your GPA after you got a PhD.
\end{itemize}
\item Dissertation/Thesis Credits
\begin{itemize}
    \item Register for CS 998 (Dissertation Proposal) for ..  
    \item Register for CS 999 (Dissertation) for ...
    \item If you perform satisfactorily in your research, you will get an \textbf{IP} (in progress) grade, which then converted to S (satisfactory) at the end of the semester when you advance to candidacy (for CS 998) or graduate (for CS 999).
    \item If you do not perform satisfactorily, I will talk to you directly. 
    \item If I cannot even reach you, I will give you an \textbf{NC} (no credit) grade. This can happen, e.g., in a case I said I would be willing to give the student a try but the student never contacted me during the semester or made any progress. 
\end{itemize}
\end{itemize}

\part{Reimbursement}

\chapter{Conference Travel}


\paragraph{Meals}
We are paid per diem for meals when we travel for conferences. The per diem rates depend on the city we are visiting. This means you don't need to save or submit receipts for meals.  If you don't eat anything (please do not do this), you will still get the per diem amount.

The first day and the last day of travel, they will not pay the full per diem rate because you are not there for the full day.


\part{Equipment and Resources}
\chapter{Connecting to Wifi and VPN}
\section{Connecting to Wifi \texttt{Eduroam}}
\begin{itemize}
\item \url{https://its.gmu.edu/knowledge-base/eduroam/}
\end{itemize}

\paragraph{Linux} \url{https://its.gmu.edu/knowledge-base/how-to-set-up-eduroam-on-linux/}
\begin{itemize}
  \item Select eduroam from the list of available wireless networks.
  \item Configure the following settings:
  \begin{itemize}
    \item Wi-Fi security: WPA \& WPA2 Enterprise
    \item Authentication: Protected EAP (PEAP)
    \item Anonymous: leave this blank
    \item Domain: gmu.edu
    \item Check No CA certificate is required
    \item PEAP version: Automatic
    \item Inner Authentication: MSCHAPv2
    \item Username: Your FULL George Mason email address (netid@gmu.edu)
    \item Password: Your Patriot Pass Password
  \end{itemize}
\item Hit Connect. You should now be connected to \texttt{eduroam}.
\end{itemize}
\section{Connecting to VPN}
TBD


\chapter{Lab Servers}

For all questions regarding lab servers, discuss in \textbf{Discord \#servers channel}.
\section{Server Specs}



We have \textbf{four} servers (\hyperref[pictures]{pictures}).

\begin{enumerate}
    \item \faMugHot{} \textbf{Boba}: ideal for AI-based projects, e.g., generative LLMs
    \begin{itemize}
        \item \href{https://www.nvidia.com/en-us/products/workstations/dgx-spark/}{NVIDIA DGX Spark}
        \item CPU: NVidia GB10 Chip, 20 core Arm (10 Cortex-X925 + 10 Cortex-A725 Arm)
        \item Unified Memory/Graphic: 128GB RAM
        \item OS: NVIDIA DGX OS (based on Ubuntu LTS)
    \end{itemize}

    \item \faPizzaSlice{} \textbf{Pizza}: ideal for projects that can leverage many cores and specialized GPU.
    \begin{itemize}
        \item Custom built
        \item CPU: 32 cores, 64 threads, 3.6GHz (AMD Ryzen Threadripper PRO 5975WX)
        \item Graphic: 2x Nvidia RTX 4090 24GB RAM
        \item Memory: 128GB RAM
        \item OS: Debian Stable
    \end{itemize}
    \begin{quote}
        This system is ideal for projects that can leverage many cores and specialized GPU.
    \end{quote}

    \item \faFish{} \textbf{Sushi}
    \begin{itemize}
        \item Lenovo ThinkStation P620
        \item CPU: 64 cores, 128 threads, 4.3GHz (AMD Ryzen Threadripper PRO 3995WX)
        \item Graphic: Nvidia RTX 5090 RTX 24GB RAM
        \item Memory: 128GB RAM
        \item OS: Debian Stable
    \end{itemize}


    \item \faUtensils{} \textbf{Taco}: ideal for experiments that need fast CPU (but not lots of core) and fast GPU.
    \begin{itemize}
        \item Alienware Aurora
        \item CPU: 16 cores, 24 threads, 5.20GHz (Intel 12th Gen Core i9-12900KF)
        \item Graphic: Nvidia RTX 3080 Ti 12GB RAM
        \item Memory: 64GB RAM
        \item OS: Debian Stable
    \end{itemize}
\end{enumerate}

\begin{warningbox}
    For Debian machines, \textbf{DO NOT} install anything other than Debian stable packages to these computers (e.g., \textbf{do not mix unstable or new repositories} to \texttt{apt} sources list). You can download and install anything in your home (local) directory. If in doubt, contact us.
\end{warningbox}

\section{Getting an Account}

To get an account on the server, send a request and your \textbf{SSH public key} (for \emph{password-less} log in).

\section{Miscs}

\begin{enumerate}
    \item You can use an IDE such as \texttt{VSCode} to do remote development on these servers (just connect to the server via the builtin remote login method in VSCode). You can also use other editors that support remote development (e.g., \texttt{Emac} + \texttt{TRAMP}).
    \item Do take advantage of the multicores these machines have. Most programming languages support multiprocessing (e.g., \texttt{Python}) and can significantly reduce program runtime.
    \item Periodically check \texttt{top} and \texttt{htop} to make sure \textbf{you} are not taking too much CPU or memory.
    \item Periodically check \texttt{/var/tmp} or \texttt{/tmp/} to make sure \textbf{you} are not taking too much space there.
\end{enumerate}

\section{Pictures}
\label{pictures}

\begin{figure}[h!]
    \centering
    \includegraphics[width=0.8\textwidth]{files/sushi.jpg}
    \caption{\faFish{} \textbf{Sushi}}
\end{figure}

\begin{figure}[h!]
    \centering
    \includegraphics[width=0.8\textwidth]{files/sushi-fan.jpg}
    \caption{\faFish{} \textbf{Sushi Cooling System}}
\end{figure}

\begin{figure}[h!]
    \centering
    \includegraphics[width=0.8\textwidth]{files/pizzabox.jpg}
    \caption{\faPizzaSlice{} \textbf{Pizza}}
\end{figure}

\begin{figure}[h!]
    \centering
    \includegraphics[width=0.8\textwidth]{files/boba.jpg}
    \caption{\faUtensils{} \textbf{Boba}}
\end{figure}




\appendix
\chapter{LoRs}

\section{Asking}

\begin{itemize}
    \item If you want a letter from me, you should \textbf{ask} at least 2 weeks in advance (ideally a month).
    \begin{itemize}
        \item \textbf{DO NOT} put my name in the application system without explicitly asking me first.
    \end{itemize}

    \item If I don't know you well enough, I will let you know and advise you to find someone else. But if you insist, I will still write you a letter.
    \item You need to \textbf{waive your right} to view the letter. It is a confidential evaluation, and I will not write one if you insist on seeing what I write.
    \item You should provide me your CV and, if you'd like, your SOP, so I can determine whether I could complement some of the things you've said.
    \begin{itemize}
        \item Feel free to provide me with any other information that you think is relevant, e.g., your grades, projects, research experience, awards, etc. I might not use them if I don't think they are relevant, but they might help if I miss something.
    \end{itemize}
\end{itemize}

\section{Writing}

\begin{itemize}
    \item I \textbf{will not ask you to draft a letter} for me. My reputation is on the line, and I do not want a student (or AI) writing a letter on my behalf.
    \item I \textbf{do not write negative letters}, i.e., I don't say bad things about you (even if you're bad). However, if I don't know you well then a neutral, short, or generic letter will not help your case given the competitiveness.
    \item It takes me about \textbf{an hour} for a strong letter for someone I know well, and about 10 minutes for a letter for someone I don't know well.
    \item In all letters, at the end I will include a short paragraph about \textbf{an area you need to improve}, regardless of whether you're super strong or weak. But I will phrase it in such a way that the graduate study environment or a good advisor can help you overcome it. It makes the letter more complete and not just full of praise.
    \item I \textbf{do not customize} the letter for specific schools.
    \begin{itemize}
        \item You customize your SOP to explain why you want to apply to school \texttt{X} and work with professor \texttt{Y}. I do not need to explain why you want to go to school \texttt{X} or why you would be a good fit for professor \texttt{Y}. Note that if I know you and \texttt{Y} very well, then I might send a letter directly to \texttt{Y} to mention you.
    \end{itemize}
    \item My letter will have the university logo and my (digital) signature.
    \begin{itemize}
        \item However, I should note that when I read a LoR, I do not pay attention to whether it has a logo or signature.
    \end{itemize}
\end{itemize}

\section{Sending}

\begin{itemize}
    \item After you put my email into the application system, it will send me a request email with a unique URL to go to. It will also give me some deadline that is \emph{likely} not the same as yours---usually later.
    \begin{itemize}
        \item In some cases the systems don't even give me a deadline (various LoR request examples are shown in my \texttt{Demystify} book, e.g., Section 3.2.3: \href{https://github.com/dynaroars/phd-cs-us/}{https://github.com/dynaroars/phd-cs-us/}).
    \end{itemize}
    \item I \textbf{do not mind if the student is anxious and sends me multiple reminders}. I will likely not respond to them. But I am not bothered by reminders. Over-communication is better than under.
    \item Most systems simply ask me to upload the letter---though a few have short questions like comparing the applicant to undergrads or grad students. As an adcom reviewer, I don't really pay attention to these comparisons---only the letter matters.
    \item I usually send my letter (in PDF) in \textbf{batch mode}, e.g., I just sent out 10+ letters all at once and cleared my inbox of these requests.
    \begin{itemize}
        \item For each request, it took me less than 30 seconds to upload the PDF and hit \texttt{submit}. So do not worry that it is time-consuming or a burden; it's not.
    \end{itemize}
    \item After I send my batch, I'll let you know I just sent everything and ask if I missed any.
    \begin{itemize}
        \item I would appreciate if you share an online spreadsheet showing what schools you have asked me to write letters for so I know what to expect (and if I miss something).
    \end{itemize}
\end{itemize}

\section{Updating}

\begin{itemize}
    \item Let me know your progress and especially your outcomes, e.g., interviews, offers, and where you eventually go to. I'd love to hear these updates from you.  
    \item This is a common courtesy. Though I probably won't remember or expect you to do this.  However, if you do this and then years later reach out for another letter, I might not be able to help you.    
\end{itemize}
\chapter{Templates}


%\bibliographystyle{abbrv}
%\bibliography{demystify.bib}

\end{document}





# Writing Papers
## LaTeX Writing 
- Use `~\cite{..}` and `~\ref{...}`: this ensures that the reference and figure numbers are not separated by a line break.
- Use `\autoref` for figures, tables, sections, equations, etc.  This will automatically add the correct prefix (e.g., "Fig.", "Table", "Section", "Eq.") to the reference. E.g., `In~\autoref{fig:ex}, ...` gives `In Fig. 1, ...`








# Managing Git repo and File structures

### Seperate Github Repo for Each Project 
  
  - Name the repo `paper_projectname`, e.g., `paper_gentree` where `gentree` is the name of the work.
  - I do not put conference name as part of the repo name because it might end up not being in that conference. Instead I use `filename` and (optionally)`git tag` ([shown below](#after-submission)) to identify submission to conference.

### File structures

1 directory per paper (which is a clone of a Github repo), e.g., `paper_gentree`
 - Within the directory, I use very few files (`paper.tex`) and no subdirectories (or at most one for figures).
    - `paper.tex`: I use a **single** `TeX` file for the entire paper.
      > Others might prefer using multiple files / directories (e.g., `intro.tex`, `eval.txt`, `sections/related.txt`). But I find it easier to just use a single file because I can easily search and replace for things and I don't have to switch between files.
    -  Even when sharing or collaborating with others, in which conflict edits can arise, a single file still works well because `Git` is pretty good at resolving conflict issues.  I have been using this approach for the last 10+ years and rarely have unresolved conflicts. Note that I do have conflicts with this approach, but they also happen with multiple files.
    - `paper.bib`: I use a single `bib` file for bibs. My collaborators sometimes put in their own bib files.
    - `arch.(pdf|png)` (optional): a diagram describing the architecture of the framework or tool

## After submission

After submitting a paper, I save a copy of the submitted pdf file and create a tag for the latest commit to keep a history of that submission.
1. Save the submitted pdf file as `VenueYearX.pdf`, where `X` is submit for the original submission version, `final` for final (camera-ready) version and `rI` for the `ith` revision for additional revisions between the original submission and final (e.g., for jounal).

    ````
    git add icse2021submit.pdf  
    ````

2. Then create an annotated tag for the commit
    ````
    git tag -a icse2021submit -m "ICSE 2021 original submission" commit_hashid (optional)
    git push origin icse2021submit
    git show icse2020submit
    ````

### After rebuttal

After submitting a rebuttal, save a copy of the reviews and rebuttal as a plain `text` file

````
git add pldi2023-reviews.txt
git commit -am "reviews and rebuttal"
````

### Miscs
- Prefer high-level description than technical one
  - People will not be interested in reading the rest if they can't understand (or get excited about) Intro
  - Start technical sections with something high level, e.g., an overview picture or algorithm.  Don't go directly into formulae and theorems.  
  - Before (or after) presenting some technical things (e.g., a theorem), explain why it is needed or useful.
- Think of a presentation as an introduction to a paper (i.e., after seeing the slides, people will want to read the paper); and think of the paper as an introduction to the work/code (i.e.., after reading the paper, people will want to use your work/code).



----


# Lab Manual





## Publications

### Authorship
 

## Rebuttal


## Data Availability
- We will *always* make the data available. So make sure prototypes, tools, results, etc are avail on Github.
- Github repo:
    - README.md: include a description of the project, how to install and use the tool.
    - For tools submitted to conferences, include scripts to generate the data and run the experiments.







# Building a machine for Neural Network Verification

We have just built a new machine **"PizzaBox"** for our research group.  Linhan and Hai built it on Oct 31st (Halloween) and it is up running since Nov 1.  The machine will be mainly used for our SE/PL research. In particular, we will use it for **neural network verification**, which is a bit different than an ML machine that focuses on training and therefore requires different hardware components and setup.

This machine is both cheaper than a machine with similar specs from Dell or Lenovo. It also runs a lot faster than computer servers available at GMU that we have access to.  We believe such customized computers are best for our purpose and hope that this post will be useful for others who are interested in building a similar machine.

## The Specs and Cost 

| Hardware | Details     | Cost | Notes |
|----------|-------------|------|-------|
| Motherboard| [ASUS Pro WS WRX80E-SAGE SE WiFi II AMD](https://www.amazon.com/gp/product/B0BZT9NF57/ref=ppx_yo_dt_b_asin_title_o05_s03?ie=UTF8&psc=1) | $1,070.52 | | 
| CPU      | AMD Ryzen™ Threadripper™ PRO 5975WX, 32-core, 64-Thread|  $2,749.98 |  from Microcenter
| GPU      | [GIGABYTE GeForce Nvidia 4090 RTX 24 GB VRAM](https://www.amazon.com/gp/product/B0BGP8FGNZ/ref=ppx_yo_dt_b_asin_title_o04_s00?ie=UTF8&th=1) |$1,699.99 | |
| RAM     | [Corsair Vengeance LPX 128GB](https://www.amazon.com/gp/product/B085WQXKM2/ref=ppx_yo_dt_b_asin_title_o05_s02?ie=UTF8&th=1) |$259.99 |
| HD  | [Western Digital BLACK 2TB SN850X NVMe](https://www.amazon.com/gp/product/B0B7CMZ3QH/ref=ppx_yo_dt_b_asin_title_o07_s00?ie=UTF8&th=1) |$119.99 | 
| PSU | [Corsair RM1000E](https://www.amazon.com/gp/product/B0BYQHWJXC/ref=ppx_yo_dt_b_asin_title_o05_s00?ie=UTF8&psc=1)|$179.99 | 
| Case | [CORSAIR 7000D AIRFLOW Full-Tower--WHITE](https://www.amazon.com/gp/product/B09444VWX2/ref=ppx_yo_dt_b_asin_title_o02_s00?ie=UTF8&th=1) | $244.99 | 
| Fan (CPU) | [Cooler Master](https://www.amazon.com/gp/product/B07H25DZ3M/ref=ppx_yo_dt_b_asin_title_o06_s00?ie=UTF8&psc=1) | 119.00 | 
| Fan (Case) | 4 x Noctua 140mm  | $23.95 | | 
|    |         |**$6,468.4**   |    |


The CPU and GPU are the two main components we need for our DNN verification. The other components are flexible.  For CPU, we go with the AMD Threadripper 32 cores (64 threads).  This is the 2nd, not top, of the line because the Threadripper 64 cores costs quite a bit and would require additional paperwork to purchase.  For GPU, we go with the current top of the line for gaming, the NVIDIA RTX 4090.  


## Timeline
- 9/11: order components
- 10/31: complete the build (1 day)

## Building It
> *Before you start*
> - make sure to discharge yourself before touching the parts.


- **Open the boxes!**
- **Install SSD, CPU, RAM:**
  - _It's often a good idea to install smaller parts onto the motherboard first._ Working in the chassis can be more challenging due to limited space.
    - **CPU:**
      - Our CPU (59775wx) uses an sWRX8 socket. The installation is a bit different from the other consumer CPUs we've seen, so it took us a while to read the manual and figure out how.
    - **SSD:**
      - Very typical NVMe SSD installation. We removed a cover from the motherboard to expose the M.2 slot, and installed the SSD in the M.2_1 slot (M.2_1 is usually directly connected to the CPU, which gives better performance).
    - **Install RAM:**
      - Nothing special about the RAM installation, either. The CPU & MB together have 8 RAM slots, each with a dedicated channel, so we used every other slot for better thermal performance. On some MBs, multiple slots might share a channel; when this happens, the MB's manual will tell you what to do.
- **Install MB:**
  - Now we can put the MB into the chassis. The chassis comes with extra standoffs and panels for ATX MBs, so we had to remove them before proceeding (because we're using an Extended-ATX, which is larger). Fit the I/O panel into the slot on the back of the chassis, and align the standoffs to the holes on the MB. Lastly, tighten the screws in an alternating pattern (so the force is evenly distributed).
- **Install PSU:**
  - Our MB was interesting; it has one MB power connector, two CPU power connectors, and one extra PCIe power connector (which is not common). The GPU uses one of the 12VHPWR connectors. We attached the cables to the PSU first, because it will be much harder to add power connectors once the PSU is in the chassis.
- **Install fans/airflow configuration:**
  - We used the conventional airflow configuration, front-to-back/top. The 140 mm fans by Noctua are really powerful yet quiet. When shopping for fans, we aimed for airflow, not air pressure, because we're not using an AIO cooler (AIO needs higher pressure to push the air through the radiators). Once we installed the fans, we routed the wires to the back of the MB, preparing for cable management.
- **Install CPU cooler:**
  - Our cooler comes with thermal paste pre-applied, so all we did was put it on the CPU and tighten the screws in an alternating pattern. This fancy cooler has a 4-pin fan connector and another SATA power connector for RGB (which none of us care about).
- **Install GPU:**
  - To install the GPU, we had to take off the PCIe slot covers. Our GPU uses two slots, so we took off two of the covers, and then we just snapped the GPU into the 1st PCIe slot. Because the 4090 is so heavy, we had to order an additional GPU supporting bracket, or the MB would be at risk.
- **Connecting cables/cable management:**
  - The last step is to connect the cables. It's simple but kind of tedious. The only potentially challenging part is connecting the system panel connectors (power button, restart button, etc.). Cable management is optional but recommended; it requires patience and some skill. 
- **Test build/install OS:**
  - After connecting everything, we pluged in the power cord and gave the system a test drive. The hardware works as expected and everything was detected by the MB. So, we continued to install the OS. However, this didn't work as well as we expected. The system boots into the installation media with no issue, but as soon as we start the installation, the keyboard stops working. After spending a couple of hours looking for a solution online, we found out that disabling SR-IOV does the trick. Strange. [link](https://www.overclock.net/threads/asus-pro-ws-wrx80e-sage-se-wifi-installing-linux.1804562/). 


## Pros and Cons
- **Pros**
    - Fun
    - Easier to upgrade (comparing PCs that uses OEM parts)
    - Full control over your machine
- **Cons**
    - unexpected problems
    - time
    - risk of braking parts


## Benchmarks


## Thoughts



---


#+TITLE: Writing
#+date: <2021-02-01 03:39>
#+description: how I write LaTeX stuff
#+filetags: computer setup blog writing
#+HTML_HEAD: <link rel="stylesheet" href="https://dynaroars.github.io/files/org.css">

Managing my $\LaTeX$ writings. 

* Writing Technical Papers
  :PROPERTIES:
  :CUSTOM_ID: writing-papers
  :END:


** Order

   I often write papers in the following order. For a concrete example, see the [[https://dynaroars.github.io/pubs/ishimwe2021dynaplex.pdf][Dynaplex paper]] published at OOPSLA'21.
   
   1. =Overview= section (Section 2).
      - start with this section, which should contain an =illustrative example=.
        - after reading this illustrative example (and perhaps the =Introduction=, the readers should have a good understanding of the approach.
        - /The goal is to have the readers (or reviewers) make up their mind whether they like or do not like the proposed work after reading this part/.
      - A good place to add some overview =figure= highlighting the main components of the technique or tool.
      - Can also summarize background or important concepts in this section
      - Typically contains:
        - An architecture or overview figure of the approach (e.g., Fig 1 in the Dynaplex paper)
        - A motivation or illustration example
          - Sec 2.1: Show how Dynaplex works at high level on an input program, what results does it produce.
        - Main definitions and background (optional)
          - Sec 2.2: general background on dynamic inference and specific features/properties of Dynaplex
   
   - =Experimental= section (Section 4).
     1. Some implementation details     
     1. Benchmark programs: descript benchmarks used
     1. Experiment Setup
        - Machine specs
        - How experiments were run (multiple times and report the mean, etc)
     1. Research Questions
       - E.g., How accurate is the technique? How efficient ?  Comparing to existing approaches
       - Experimental results are designed to answer these questions
       
   - =Algorithm/Techniques= (Section 3)
     
   - =Related Work= (Section 5)
     - Break into multiple subsections if necessary
     - After each subsection that talks about related work, have a new paragraph and talk about why my work is different/related to other work.
       
   - =Conclusion= (Section 6).
     - This should be the /conclusion/ not the /summary/.
     - Future work can also go here
       
   - =Introduction= (Section 1).
     - write this last, because by the time I write the Intro, I would already have all needed info (e.g., so I can summarize the diff with state of the art, technical contributions, and results.
     - Typically contains several main things:
       - What is the problem?  Why is it important (i.e., if we can solve it, then what do we accomplishes)? 
       - Existing work:  basically answering why it is challenging (i.e., why existing work is not sufficient)?
       - The proposed approach of this paper: how does it address the limitations of existing work listed above?  What other benefits does it have?
       - Results: briefly summarize the results.
   
   

** Paper naming convention
   :PROPERTIES:
   :CUSTOM_ID: paper-naming-convention
   :END:

- I use Google Scholar's convention  =lastnameoffirstauthor~publishedyear~firsttitleword=, e.g., =nguyen2021gentree.pdf= for the paper  =KimHao Nguyen and ThanhVu Nguyen. GenTree: Using Decision Trees to Learn Interactions for Configurable Software=
- I also use this convention for =bibtex= entry and reference, e.g.,
  =GenTree paper~\cite{nguyen2021gentree}=
- If there's a conflict, then append the name with a number, e.g.,
  =nguyen2021gentree1=, =nguyen2021gentree2=


  
** Citation
   :PROPERTIES:
   :CUSTOM_ID: citation
   :END:
   
- As with [[#paper-naming-convention][paper naming]], for citation I follow Google Scholar's convention and thus often use Google Scholar to search for  =bibtex= entry of a paper I want to cite.  For example, to cite the GenTree example paper, I search for it on Google Scholar and get the =bibtex= content.
  #+begin_src bibtex
    @inproceedings{nguyen2021gentree,
      title={GenTree: Using decision trees to learn interactions for configurable software},
      author={Nguyen, KimHao and Nguyen, ThanhVu},
      booktitle={2021 IEEE/ACM 43rd International Conference on Software Engineering (ICSE)},
      pages={1598--1609},
      year={2021},
      organization={IEEE}
    }
  #+end_src
Then I copy/paste this =bibtex= content to my the bib file of my project, and cite the paper with =\cite{nguyen2021gentree}=.
  
** Do and Don't
   - Do
     - Split into multiple paragraphs.  Long paragraphs are simply hard to read!

   - Do not

     - Delay talking about the main work until the 2nd page of the introduction.  It's not good if people read a whole page and still don't know what you're proposing!
       
     - Use unnecessary roadmaps.  Eg.., at the end of an intro, many papers have something like this
       #+begin_src txt
       This paper is organized as follows. In Section 2, we describe the background. In Section 3, we establish theorem X. In Section 4, we describe algorithm Y and implementation Z. In Section 5, we present experimental results from our implementation and compare to the state of the art.  In Section 6, we review related work and conclude in Section 7.
       #+end_src
       - This might make sense for a long Ph.D. thesis or book (in fact, I wouldn't use this roadmap even for these), but doesn't make much sense for a 10 page paper where one easily guess all these from the section names?

     - Similarly, we can elimiate roadmap at the beginning of a section
        #+begin_src txt
        In this section, we first describe X, then present Y.  We also do Z.
        #+end_src
        - This might be OK, especially for a long section.  But in general I prefer just start talking about the main stuff in the section.

          
   
** Trimming / Space saving
   - Rewrite and shorten sentences!! 
     

** Miscs
- One setence per line
  
  #+begin_example
      This is a line.
      This is another, longer line.
  #+end_example

- If possible, put figures, tables, code, etc on top, e.g., using
  =\begin{figure}[t]=
- Rewrite to avoid lines with single or lone words

  #+begin_example
      It just looks
      bad.
  #+end_example

- Comment:  so that my coauthors and I can add comments to paper, e.g., =\tvn{This is my comment}= 
  #+begin_src latex
    \newcommand{\mycomment}[3][\color{red}]{{#1{[{#2}: {#3}]}}}
    \newcommand{\tvn}[1]{\mycomment[\color{red}]{TVN}{#1}}
   #+end_src


* Managing Git repo and File structures
   :PROPERTIES:
   :CUSTOM_ID: git-repo
   :END:

- 1 Git repo for each paper
  - Name the repo =paper_name=, e.g., =paper_gentree= where =gentree= is the name of the work. 
  - I do not put conference name as part of the repo name because it might end up not being in that conference. Instead I use git =tag= ([[#after-submission][shown below]]) to identify submission to conference.

** Paper structure
    :PROPERTIES:
    :CUSTOM_ID: paper-structure
    :END:

- 1 directory per paper (which is a clone of a Github repo as described [[#git-repo][above]]), e.g., =paper_gentree/=
- Within the directory, I use very few files:
  - =paper.tex=: I use a single $\TeX$ file for the entire paper.
    - Others like to split into multiple files (e.g., =intro.tex=, =eval.txt=, =related.txt=, etc). But I find it easier for me to just use 1 file. Even when sharing or collaborating with others, in which conflict edits can arise, a single file still works well as git is pretty good at resolving conflict issues.
  - =paper.bib=: I use a single =bib= file for bibs. My collaborators sometimes put in their own =bib= files. 
  - =arch.(pdf|png)= (optional): a diagram describing the architecture of the framework or tool

** After submission
   :PROPERTIES:
   :CUSTOM_ID: after-submission
   :END:

After submitting a paper, I save a copy of the submitted pdf file and create a tag for the latest commit to keep a history of that submission.
   
1. Save the submitted pdf file as =VenueYearX.pdf=, where ~X~ is ~submit~ for the original submission version, ~final~ for final (camera-ready) version, and ~rI~ for the $i^{th}$ revision for additional revision submissions between the original submission and final (e.g., for jounal).
   #+begin_example
     git add icse2021submit.pdf  
   #+end_example
1. Create an annotated =tag= for the commit
  #+begin_example bash
    git tag -a icse2021submit -m "ICSE 2021 original submission" commit_hashid (optional)
    git push origin icse2021submit
    git show icse2020submit
  #+end_example

** After rebuttal
After submitting a rebuttal, save a copy of the reviews and rebuttal as a plain ~text~ file
   #+begin_example
     git add pldi2023-reviews.txt
     git commit -am "reviews and rebuttal"
   #+end_example


* Quick Fixes
   1. Avoid passive tense
   - Know the difference between /Such as/  and /like/
   - Avoid long paragraphs
   - Spell out numbers less than 10 (/three/ steps instead of /3/ steps).
   - Use Eq. X instead of Equation X
   - Don't put /etc/ in e.g., or such as
   - In order to do this -> To do this
   - New paragraph for new thought/idea. Sentences in the same paragraph should be connected.

* Presentation/Talk
  :PROPERTIES:
  :CUSTOM_ID: presentation
  :END:

- Slides

  - Put numbers on every slide: so that people can refer to slide X
  - For an N-minute presentation: create about N slides
  - Don't put too much words, also no need to put full sentences





   
* Useful Links
  - http://www.cs.cmu.edu/%7Ejrs/sins.html



  ----


#+TITLE: GMU CS FAQ
#+date: <2022-09-05 Mon>
#+description: Things about GMU CS
#+filetags: gmu cs bsd guide 

#+HTML_HEAD: <link rel="stylesheet" href="https://dynaroars.github.io/files/org.css">



* Funding
** GRA Offerring
*** For grant fund GRAs (e.g., from NSF)
   - Contact Terri and her team to make such offers and to generate contracts
*** For CEC fund GRAs (e.g., from your startup GRAs)
   - Contact Jennifer Lee and her team to make such offer and to generate contracts

** Unrestricted Gift / Industrial Awards
This will go through *[[https://gmuf.org/][GMU Foundation (GMUF)]]*.  GMU will create an account for you for your unrestricted gift. You can use the same account for multiple awards (e.g., convenient because you don't have to remember too many account/org #'s)
** Contact
   - *Beth Cantrell* ecantrel@gmuf.org (Vice President and Chief Financial Officer of GMUF)
   - *Kira Woitek* kwoitek@gmu.edu (Director of Development, CEC)

  Beth is the financial point of contact at GMUF.  Kira's office helps with creating an account for the gift.

  
*** GMU Charge
GMUF assesses a 6% fee on any current-use (spendable) gifts at the time of the gift.  So if you have a 30K gift, then GMU takes $1800. The fee supports future Advancement initiatives as well as the Foundation's infrastructure necessary to administer funds.

* Purchase
* Teaching
* Travel
** What can and cannot be reimbursed for travelling
*** Can be purchased through P-card
    - Airline tickets
    - Conference registration.
*** Can be reimbursed (purchase individually and get reimbursed)
  - Everything else involved in travel must be purchased by the individual and reimbursed.  This includes things like the hotel, Uber, taxi's, food, car rentals, subways,  parking, tolls and more.
*** Cannot be reimbursed
- *Insurance associated with rental cars in the United States* (see car rental information below)
- Damage to personal vehicles
- Expenses related to personal negligence of the traveler, such as fines
- Services to gain entry to locked vehicles
- Towing charges
- Luxury vehicle rentals
- Excessive parking and gasoline charges
- University parking fees
- Local commuting expenses
- Travel between lodging and places where meals may be taken

** Per diem
- https://fiscal.gmu.edu/travel/learn-about-meals-per-diem-rates/

** Reimbursement
** Max meal
** Car Rental
- GMU/Commonwealth of VA has its own insurance  for rental cars in the US (i.e., you DO NOT need to get them — if you do, you are not getting reimbursed).  https://risk.gmu.edu/insurance/vehicle/
  - For International, you can get reimbursed. 
  - However, to use this insurance, you need to submit this MVR form  and schedule to take a training course (Zoom,  https://risk.gmu.edu/training/driver-training/). You only do this once.
  - GMU prefers Enterprise for car rentals — but I don’t think they would deny you the insurance if you use a different company.
* Miscs
** Creasing a website



  




# Research
> This is about publishing papers.

## Advise
- **Short/workshop papers can be very useful**
- 

# Funding
> This is about obtaining grants to support students and summer salary for faculty.




## Advise

- **Volunteer to serve in NSF Panels**
  - Contact program managers in your area and ask to volunteer to serve on their NSF Panels. In general NSF PMs are very nice and they'd love to support junior faculty.
  
- **Done with draft early to get feedback**

- **Have preliminary results**

- **Have supporting letters**: a supporting letter can help.  Just like having preliminary results, these letters show that you will have collaborators and expoert that can help with the proposed work.  Support letters can be from anyone with the right background that fit with the proposed work (e.g., a high school teacher for hs outreach activity, a collaborator from the industry).

# Mentoring
> This is about working with and mentoring students.

# Teaching
> This is about teaching graduate and undergraduate courses

# Service
> This is about serving in your professional communities (e.g., serving in program committees of conferences) and department/universities (e.g., serving in the Ph.D. admission committee)





# Research
> This is about publishing papers.

## Advise
- **Short/workshop papers can be very useful**
- 

# Funding
> This is about obtaining grants to support students and summer salary for faculty.

## Types of Funding
### Grants
  - Freedom to pursue basic research

#### Where to get them?  
  - NSF
    - Has REU supplement (e.g., 16K for each NSF grant)
    - Typically has a 1 year no-cost extension
  - DARPA
  - Army/Navy/Airforce

### Contract
  - Usually have some sort of "deliverables" or report requirements
  - Doesn't give lots of freedom as other grants.

### Internal Funding
  - Start up package
    - when you join the faculty at a school (either startiong completely new assistant profs. or moving from another place)
    - Typically has GRA and Summer supports. Also funding for various other things such as equipments, conference travels. 
    - Expire after 3 years or so
    
  - University funding, e.g., "Seed" funding
    - Usually small, 20K or below, 1 or 2 years.
    - Might not support faculty summer salary
    - The goal is to help you start some project to get external funding later
    
  - Statewide Initiative / Research Center
    - Similar to University funding but usually give more, e.g., 50 K or so
    - Goal is to initiate a project or collaboration (e.g., with faculty at another school in the state) to prepare for larger external proposal
  
### Industrial/Foundation Funding  
  - Unrestricted gifts: The *best* kind of funding
    - No indirect cost: university cannot charge indirect cost to these funds
    - No expiration: you can just leave them there and use in case of emergency only
    - Can be used for much broader range of things: e.g., paying for meals for your research group
    - Can be highly competitive:  you're competing against everyone

#### Where to get them?
    - Industrial: Amazon, Google, Facebook, Microsoft, etc
    - Foundation: Sloan, 

## Advise

- **Volunteer to serve in NSF Panels**
  - Contact program managers in your area and ask to volunteer to serve on their NSF Panels. In general NSF PMs are very nice and they'd love to support junior faculty.
  
- **Done with draft early to get feedback**

- **Have preliminary results**

- **Have supporting letters**


# Mentoring
> This is about working with and mentoring students.

# Teaching
> This is about teaching graduate and undergraduate courses

# Service
> This is about serving in your professional communities (e.g., serving in program committees of conferences) and department/universities (e.g., serving in the Ph.D. admission committee)


# Links
- https://vijay03.github.io/asstprofbook/
